\documentclass[11pt]{article}
\usepackage{fontspec}
\usepackage[french]{babel}
\usepackage{amsmath}
\usepackage{amssymb}
\usepackage{amsthm}
\usepackage{mathrsfs}
\usepackage{tikz-cd}
\usepackage{enumitem}

\newcommand{\printedpage}[1]{\marginpar{\footnotesize p.~#1}}
\newcommand{\msp}[1]{\marginpar{\footnotesize [#1]}}
% hooked arrow with optional label above
\newcommand{\hto}[1][]{\overset{#1}{\lhook\joinrel\longrightarrow}}
% circled letter
\newcommand{\circled}[1]{\tikz[baseline=(c.base)]\node[draw,circle,inner sep=1pt](c){#1};}
% integral of a diagram  #1 -> #2, #1 -> #3
\newcommand{\intdiag}[3]{\begin{tikzcd}[row sep=small,column sep=small,ampersand replacement=\&,baseline=(\tikzcdmatrixname-2-1.base)]
 \& #2 \\
 \displaystyle\int #1 \arrow[ur] \arrow[dr] \& \\
 \& #3
\end{tikzcd}}

\begin{document}

\printedpage{103}
% head: A. Grothendieck : Dérivateurs, Ch. XII, Édition des Universités de Montpellier II et Paris VII

\noindent\textbf{Corollaire 4.} (Moyennant W(7 bis) (\footnote[104]{Ici on a besoin de W(7 bis) forme forte, i.e. $i_0\in W \iff i\in W$ (cf. annotation th. 4).}).) \textit{Soit $U \hto[i] X$ une immersion ouverte telle que $T=X\backslash U$ soit isomorphe à la catégorie ponctuelle $e$. Pour que $i\in W$, il faut et il suffit que $i$ soit $W$-asphérique, i.e. que $U/t$ (où $t$ est l'unique objet de $T$) soit $W$-asphérique.}

Il suffit bien sûr de prouver que
\[
i\in W \Longrightarrow U/t \ W\text{-asphérique.}
\]

\msp{page 150}
Considérons le carré cartésien
\[
\begin{tikzcd}[row sep=large,column sep=large]
U/t \arrow[r, hook, "i/t"] \arrow[d, "p"'] & X/t \arrow[d, "q"] \\
U \arrow[r, hook, "i"] & X,
\end{tikzcd}
\]
nous allons prouver qu'il est $W$-cocartésien. Moyennant W(7 bis), il s'ensuivra que $i\in W$ implique $i/t\in W$, et comme $X/t$ est $W$-asphérique (ayant un objet final), cela signifie que $U/t$ est $W$-asphérique, ce qu'on veut prouver.

Il suffit, en vertu du corollaire 2, de prouver que le carré est cocartésien dans Cat, puisque $i/t$ est une immersion ouverte ($i$ l'étant), et $X/t\longrightarrow X$ est un isomorphisme local, donc $W$-lisse. On voit que $(X/t)\backslash(U/t)\simeq e$ (catégorie ponctuelle), d'ailleurs l'unique élément de $(X/t)\backslash(U/t)$ est l'élément final $t_0 : t\xrightarrow{\mathrm{id}_t} t$ de $X/t$. Ainsi $H_{t_0}\simeq e_{U/t}$ (préfaisceau final), et $p_{!}(H_{t_0})\simeq H'_t \overset{\text{déf}}{=} (u\longmapsto \mathrm{Hom}_X(u,t))$, ce qui caractérise la somme amalgamée dans Cat
\dots

Le corollaire 4 n'est autre que le corollaire non démontré de la proposition 7 (p. 85).

Voici une variante du corollaire 4.

\noindent\textbf{Corollaire 5.} (Moyennant W(7 bis) (\footnote[105]{forme forte.}).) \textit{Soit $i : U \hto X$ une immersion ouverte, soit $t_0\in\mathrm{Ob}\,T$ un point \textup{fermé} de $T=X\backslash U$, tel que $\{t_0\}$}
\msp{page 151}
\textit{soit isomorphe à la catégorie ponctuelle, i.e. $\mathrm{Hom}_X(t_0,t_0)=\{\mathrm{id}\}$. Supposons que pour tout $t\in T\backslash\{t_0\}$, $U/t$ soit $W$-asphérique. Pour que l'on ait $i\in W$, il faut et il suffit que $i$ soit asphérique, i.e. que $U/t_0$ soit également $W$-asphérique.}

Il faut prouver que $i\in W$ implique que $U/t_0$ est $W$-asphérique. Considérons l'inclusion
\[
i_0 : U\longrightarrow U' \overset{\text{déf}}{=} X\backslash\{t_0\},
\]
l'hypothèse sur les $U/t$, pour $t\in T\backslash\{t_0\}$, signifie que $i_0$ est $W$-asphérique. Donc $i_0$ est dans $W$, et comme le composé $i : U\xrightarrow{i_0} U'\xrightarrow{i'} X$ est dans $W$, cela implique que
% foot: Version 7 décembre 2002
\printedpage{104}
% head: A. Grothendieck : Dérivateurs, Ch. XII, Édition des Universités de Montpellier II et Paris VII
$i' : U'\longrightarrow X$ est dans $W$. Nous sommes maintenant dans la situation du corollaire 4, qui implique que $i'$ est $W$-asphérique. Donc $i=i'i_0$ est $W$-asphérique, comme composé de deux morphismes $W$-asphériques, q.e.d.

\noindent\textsc{Démonstration de th. 4.}

\noindent\circled{a} Critère de lissité de $\varphi_Q : \mathcal{X} = \intdiag{X_0}{X_1}{X_2} \longrightarrow Y$.

Comme $\{1\}$, $\{2\}$ sont des ouverts de $\Phi = \begin{tikzcd}[row sep=small,column sep=small,baseline=(\tikzcdmatrixname-2-1.base)]
 & 1 \arrow[dl] \\
 0 & \\
 & 2 \arrow[ul]
\end{tikzcd}$, $i_1 : X_1\longrightarrow\mathcal{X}$ et $i_2 : X_2\longrightarrow\mathcal{X}$ sont des immersions ouvertes. Comme $\varphi_Q\circ i_1 = f_1$,
\msp{page 152}
$\varphi_Q\circ i_2=f_2$, on voit déjà que
\[
f_1 \text{ et } f_2 \text{ lisses} \iff \varphi_Q \text{ lisse sur l'ouvert réunion de } i_1(X_1) \text{ et } i_2(X_2)\text{, i.e. sur } \mathcal{X}\backslash i_0(X_0).
\]
Il reste à exprimer la lissité de $\varphi_Q$ en les points $x_0$ de $i_0(X_0)$. Soient $x_1$, $x_2$, $y$ les images de $x_0$ dans $X_1$, $X_2$, $Y$ respectivement. Il faut exprimer que les fibres de
\[
(\varphi_Q)/x_0 : \mathcal{X}/x_0\longrightarrow Y/y \tag*{$(*)$}
\]
sont asphériques. Nous allons pour cela expliciter $\mathcal{X}/x_0$ et la flèche $(*)$ en termes du carré commutatif
\[
(**)\qquad\qquad
\begin{tikzcd}[row sep=large,column sep=large]
X_0/x_0 \arrow[r, "\alpha_1/x_0"] \arrow[d, "\alpha_2/x_0"'] & X_1/x_1 \arrow[d, "f_1/x_1"] \\
X_2/x_2 \arrow[r, "f_2/x_2"] & Y/y,
\end{tikzcd}
\]
induit par $Q$. Je dis

\noindent\textbf{Lemme 1.} \textit{Soit un carré commutatif}
\[
Q \quad = \quad
\begin{tikzcd}[row sep=large,column sep=large]
X_0 \arrow[r, "\alpha_1"] \arrow[d, "\alpha_2"'] & X_1 \arrow[d, "f_1"] \\
X_2 \arrow[r, "f_2"] & Y
\end{tikzcd}
\]
\textit{dans} Cat, \textit{$x_0\in\mathrm{Ob}\,X_0$, $x_1$, $x_2$, $y$ ses images dans $X_1$, $X_2$, $Y$ respectivement, $Q/x_0$ le carré $(**)$ ci-dessus, induit par $Q$ sur les localisées $X_0/x_0$ etc., $\mathcal{X} = \intdiag{X_0}{X_1}{X_2}$. On a alors un}
% foot: Version 7 décembre 2002
\printedpage{105}
% head: A. Grothendieck : Dérivateurs, Ch. XII, Édition des Universités de Montpellier II et Paris VII
\msp{page 153}
\textit{isomorphisme canonique}
\[
\mathcal{X}/x_0 \quad\simeq\quad \intdiag{X_0/x_0}{X_1/x_1}{X_2/x_2} ,
\]
\textit{et la flèche $(\varphi_Q)/x_0$, localisée de $\varphi_Q : \mathcal{X}\longrightarrow Y$ en $x_0$, s'identifie à la flèche $\varphi_{(Q/x_0)}$ :}
\[
\varphi_{(Q/x_0)} : \quad \intdiag{X_0/x_0}{X_1/x_1}{X_2/x_2} \qquad\longrightarrow\qquad Y/y.
\]

\noindent\textsc{Démonstration} du Lemme. On sait que
\[
\mathcal{X}/x_0\longrightarrow\mathcal{X}
\]
est Cat-fibrée (scindée) (à fibres discrètes), et comme
\[
\mathcal{X}\longrightarrow\Phi
\]
est aussi Cat-fibrée (scindée), il s'ensuit qu'il en est de même de la flèche composée
\[
\mathcal{X}/x_0\longrightarrow\Phi.
\]
Donc pour `calculer' $\mathcal{X}/x_0$, en tant que catégorie sur $\Phi$, il faut expliciter ses trois fibres, et les deux foncteurs de changement de base qui les relient. On trouve aussitôt que ces fibres sont canoniquement isomorphes à
\[
X_0/x_0, \qquad X_1/x_1, \qquad X_2/x_2,
\]
et on admettra (!) que les foncteurs changement de base
\msp{page 154}
\[
X_0/x_0\longrightarrow X_1/x_1, \qquad X_0/x_0\longrightarrow X_2/x_2
\]
(définis à isomorphisme unique près) ne sont autres que les localisés $\alpha_1/x_0$ et $\alpha_2/x_0$ de $\alpha_1 : X_0\longrightarrow X_1$ et $\alpha_2 : X_0\longrightarrow X_2$. Cela montre donc qu'on a un isomorphisme
\[
\mathcal{X}/x_0 \quad\simeq\quad \intdiag{X_0/x_0}{X_1/x_1}{X_2/x_2} ,
\]
% foot: Version 7 décembre 2002
\printedpage{106}
% head: A. Grothendieck : Dérivateurs, Ch. XII, Édition des Universités de Montpellier II et Paris VII
et l'on admettra aussi que la flèche $(\varphi_Q)/x_0 : \mathcal{X}/x_0\longrightarrow Y/y$ s'identifie à $\varphi_{(Q/x_0)}$.

Pour exprimer la $W$-lissité de $\varphi_Q$ en $x_0\in\mathrm{Ob}\,X_0\subseteq\mathrm{Ob}\,\mathcal{X}$, on est donc ramené à exprimer la $W$-asphéricité des fibres de
\[
\varphi_{(Q/x_0)} : \quad \intdiag{X_0/x_0}{X_1/x_1}{X_2/x_2} \qquad\longrightarrow\qquad Y/y.
\]
Pour expliciter ces fibres, on a besoin d'un autre lemme auxiliaire

\noindent\textbf{Lemme 2} (\footnote[106]{NB C'est le b) du th. 4.}). \textit{Soit un carré commutatif dans} Cat
\[
Q \quad = \quad
\left(
\begin{tikzcd}[row sep=large,column sep=large]
X_0 \arrow[r, "\alpha_1"] \arrow[d, "\alpha_2"'] & X_1 \arrow[d, "f_1"] \\
X_2 \arrow[r, "f_2"] & Y
\end{tikzcd}
\right),
\]
\textit{d'où}
\[
\varphi_Q : \quad \left( \intdiag{X_0}{X_1}{X_2} \right) \qquad\longrightarrow\qquad Y,
\]

\msp{page 155}
\textit{et soit $y\in\mathrm{Ob}\,Y$. Alors la fibre de $\varphi_Q$ en $y$}
\[
\mathcal{X}_y = \varphi_Q^{-1}(\ \underbrace{\{y\}}_{\text{catégorie ponctuelle}}\ )
\]
\textit{est donnée par (comme catégorie au dessus de $\Phi$)}
\[
\mathcal{X}_y \quad\simeq\quad \intdiag{X_{0\,y}}{X_{1\,y}}{X_{2\,y}} .
\]

\noindent\textsc{Démonstration.} Il est immédiat qu'on a
\[
\mathcal{X}_y\cap\underbrace{\mathcal{X}_s}_{X_s}\simeq(X_s)_y \qquad\quad \text{pour } s\in\mathrm{Ob}\,\Phi=\{0,1,2\},
\]
% foot: Version 7 décembre 2002
\printedpage{107}
% head: A. Grothendieck : Dérivateurs, Ch. XII, Édition des Universités de Montpellier II et Paris VII
et (le hasard fait bien les choses) les foncteurs changement de base
\[
X_0\longrightarrow X_1, \qquad X_0\longrightarrow X_2
\]
envoient la sous-catégorie (non pleine) $X_{0\,y}$ de $X_0$ dans les sous-catégories (non pleines) $X_{1\,y}$ et $X_{2\,y}$ de $X_1$ et $X_2$ respectivement.
\[
\begin{tikzcd}[row sep=large,column sep=large]
X_{s\,y} \arrow[r] \arrow[d] \arrow[dr, phantom, "\scriptstyle\text{cart.}"] & \mathcal{X}_y \arrow[r] \arrow[d] \arrow[dr, phantom, "\scriptstyle\text{cart.}"] & \{y\} \arrow[d] \\
X_s \arrow[r] \arrow[d] \arrow[dr, phantom, "\scriptstyle\text{cart.}"] & \mathcal{X} \arrow[r] \arrow[d] & Y \\
\{s\} \arrow[r] & \Phi &
\end{tikzcd}
\]
On a donc une sous-catégorie fibrée de $\mathcal{X}$, sous réserve de vérifier que les flèches strictement cartésiennes
\[
x_0\longrightarrow\alpha_1(x_0), \qquad x_0\longrightarrow\alpha_2(x_0)
\]
dans $\mathcal{X}$, pour $x_0$ dans $X_{0\,y}$, sont dans $\mathcal{X}_y$, i.e. sont au-dessus de la flèche identique de $y$. Or c'est là la propriété essentielle
\msp{page 156}
du foncteur $\varphi_Q : \mathcal{X}\longrightarrow Y$, qu'on a oublié d'expliciter dans 3.2. Ça a vraiment l'air de marcher (comme sur des roulettes!)

Appliquons le Lemme 2 au diagramme carré $Q/x_0$ et à $\varphi_{(Q/x_0)}$ (page 154), on voit donc que les fibres à examiner en les $u\in\mathrm{Ob}\,Y/y$ ($u : y'\longrightarrow y$) peuvent s'écrire comme les intégrales
\[
\mathcal{X}(x_0,u) \quad\overset{\text{déf}}{=}\quad \intdiag{(X_0/x_0)_u}{(X_1/x_1)_u}{(X_2/x_2)_u} .
\]
Utilisons maintenant le fait que $f_1$, $f_2$ sont $W$-lisses (exprimant que $\varphi_Q$ est $W$-lisse dans $\mathcal{X}\backslash X_0$), donc que $(X_1/x_1)_u$ et $(X_2/x_2)_u$ sont $W$-asphériques. Dans le carré $W$-cocartésien
\[
\begin{tikzcd}[row sep=large,column sep=small]
 & (X_1/x_1)_u \arrow[dr, "{i_1(x_0,u)}"] & \\
(X_0/x_0)_u \arrow[ur, "{\alpha_1(x_0,u)}"] \arrow[dr, "{\alpha_2(x_0,u)}"'] & & (\mathcal{X}/x_0)_u\simeq\mathcal{X}(x_0,u) \\
 & (X_2/x_2)_u \arrow[ur, "{i_2(x_0,u)}"'] &
\end{tikzcd}
\]
% foot: Version 7 décembre 2002

\end{document}
